\chapter{Discussion and Interpretation}\label{ch:discussion}

This chapter interprets the results. Measured observations, interpretations and untested hypotheses are kept apart
throughout.

\section{Why Table-Specialist works, and when}

\begin{observation}
\begin{itemize}
\item Fine-tuning on validated pairs raises zero-shot accuracy substantially on both checkpoints: \ZSiv{} $+0.138$ on
base and $+0.111$ on Instruct.
\item The gain concentrates on the test sets whose failures are formatting and quoting: WikiSQL, and WikiTQ to a lesser
degree.
\item With the few-shot prompt, the gain shrinks to near zero or turns negative (\cref{tab:gen-prompt}).
\end{itemize}
\end{observation}

\begin{interpretation}
Much of what self-generated training teaches an 8B model is \emph{how to answer the zero-shot prompt}: a well-formed,
quoted, single SQL block. The few-shot prompt teaches the same thing by example. The specialist's further gains are
smaller and test-set-specific. This reading fits the observation that few-shot vanilla Instruct matches zero-shot TS
(\cref{sec:ext-fewshot}). It does not mean fine-tuning is useless: the zero-shot prompt is about 2,000 tokens shorter,
and the vote gains on top of TS.
\end{interpretation}

\section{Validation, volume and rounds}

\begin{observation}
\begin{itemize}
\item NV beats TS in round~1, then falls every round.
\item TS rises on base and is flat on Instruct.
\item Validation discards 75--89\% of generated pairs (\cref{tab:gen-validation}).
\end{itemize}
\end{observation}

\begin{interpretation}
Validation trades volume for precision. The trade pays off only when a model's own errors feed back into its training,
that is, from round~2 onward. That is exactly when NV declines. In round~1, both arms' data come from the same vanilla
model, and NV's sixfold volume wins.
\end{interpretation}

\begin{hypothesis}
Instruct does not gain from rounds because its round-1 adapter already captures most of what this corpus can teach
through about 500 validated pairs. Restarting from vanilla means a later round's model sees no more data than an
earlier one, only data written by a better generator. Testing this needs either cumulative data across rounds or
continued training. Neither was run.
\end{hypothesis}

\section{When does repairing rejected data help?}

The central new finding is that repair-TS's value depends on the setting:

{\small
\rowcolors{2}{rowshade}{white}
\begin{longtable}{@{}p{3.0cm}p{2.1cm}p{2.4cm}p{3.3cm}p{3.6cm}@{}}
\caption{Repair-TS across settings, at round~3 unless stated.}\label{tab:repair-summary}\\
\toprule\rowcolor{headshade}Setting & Generator & Repairer & Effect & Recovery\\\midrule\endfirsthead
\toprule\rowcolor{headshade}Setting & Generator & Repairer & Effect & Recovery\\\midrule\endhead
\bottomrule\endfoot
NL-to-SQL, base & Base & Instruct (stronger) & $+0.0421$, $p<0.0001$ & +52\% (round 1), about +20\% later\\
NL-to-SQL, Instruct & Instruct & Instruct (same) & $-0.0014$, $p=0.55$ & +15--23\%\\
NL-to-Scala, Instruct & Instruct & Instruct (same) & $+0.0167$, $p<0.0001$ & the same pairs\\
NL-to-R, Instruct & Instruct & Instruct (same) & $+0.0244$, $p<0.0001$ & +25--31\% (round 1), +16--23\% later\\
Error detection, Instruct & Vanilla Instruct & Vanilla Instruct & $+0.011$, $p=0.58$ (round~1: $+0.027$, $p=0.03$) &
about 490--700 errors per round\\
\end{longtable}
}

\begin{interpretation}
Two readings fit the SQL rows; our experiments do not separate them.
\begin{enumerate}
\item \textbf{Repair as distillation.} On base, the repairer is a stronger model than the generator, and the repaired
programs carry its knowledge. On Instruct, generator and repairer are the same model, so repaired programs add nothing
new. Recovered pairs then concentrate on easy failure types (\texttt{disagree}, not \texttt{always\_empty}), which is
why they are worth less than fresh first-pass pairs.
\item \textbf{Repair as recovery of the weak side.} Repair helps where first-pass data is the bottleneck. That is base
SQL (14.7\% kept) and Scala on Instruct (vanilla \ZSiv{} 0.19), but not Instruct SQL, which is already strong and keeps
enough pairs.
\end{enumerate}
The Scala row favours reading~2. Repair helps Scala on Instruct with the same model as repairer, which reading~1 alone
does not predict. It is also consistent with reading~1 to the extent that the Instruct \emph{repairer}, prompted with
execution feedback, writes better Scala than the Instruct \emph{generator} does on the first pass.
\end{interpretation}

\begin{hypothesis}
The value of recovered data is set by the gap between the repairer's ability and the trained model's ability, in the
language being trained. Testing this would need:
\begin{itemize}
\item base repairing itself with a prompt that works;
\item Instruct writing fresh pairs for base, a distillation control without repair;
\item a weaker and a stronger model of another family.
\end{itemize}
\end{hypothesis}

\section{Why DPO failed while repair succeeded on base}

Both use the rejected data, but in opposite ways. DPO lowers the probability of plausible-but-wrong programs. With
LoRA at $2\times10^{-5}$, it moved the policy far from the reference: reward margins of 3--5, and execution errors up
to fourfold. Repair turns the same rejected data into new \emph{positive} examples that pass the unchanged validator.
The evidence favours using the validator to \emph{select} better data rather than to define negatives. Whether a better
preference method (another loss, KL anchoring, or many more pairs) would change this is open.

\section{Sensitivity to seeds and stability across rounds}

\begin{itemize}
\item \textbf{Seed spread.} Spreads of \Mviii{} reach 0.02 for TS. That is as large as several method differences we
examined.
\item \textbf{What is robust.} Every claim in this report that we call consistent has the same sign on all three seeds.
Claims with mixed seeds are reported as such (bestof4, Instruct repair-TS, the same-data control's volume effect on
base).
\item \textbf{Stability.} Repair-TS is markedly more stable across seeds than TS on both tasks (base SD 0.0051 against
0.0210; Scala \ZSiv{} SD 0.0031 against 0.0129).
\end{itemize}

\begin{hypothesis}
Recovered data makes each seed's training set larger and more representative, which damps seed-specific quirks of
small first-pass sets.
\end{hypothesis}

\section{Reproducibility considerations}

\begin{itemize}
\item \textbf{Within this project}, everything is reproducible from saved artefacts:
  \begin{itemize}
  \item every generation is cached;
  \item every adapter and per-case verdict is saved;
  \item every table in this report is generated from result files;
  \item the cross-check in \cref{app:crosscheck} confirms 24 of 24 headline values.
  \end{itemize}
\item \textbf{Across projects}, three factors limit exact reproduction:
  \begin{itemize}
  \item GitTables sampling depends on our filters, which are fully specified;
  \item vLLM's batched greedy decoding is not guaranteed bit-identical across runs;
  \item the model weights are pinned by revision.
  \end{itemize}
\item \textbf{The paper itself} cannot be reproduced exactly by anyone without its corpus, case lists and settings.
\end{itemize}

\section{What is established and what remains uncertain}

\begin{keypoint}
\textbf{Established} (consistent across seeds, $p<0.0001$, controls where relevant):
\begin{itemize}
\item Instruct TS approximately reproduces the paper's NL-to-SQL results.
\item TS beats vanilla and NV at round~3.
\item NV declines over rounds.
\item Repair-TS helps base NL-to-SQL and Instruct NL-to-Scala, and does not help Instruct NL-to-SQL.
\item On base, repaired data is worth more than as many first-pass pairs.
\item DPO on validator labels harms.
\item Execution voting helps every model.
\end{itemize}
\textbf{Uncertain:}
\begin{itemize}
\item why base does not reproduce the vanilla scores;
\item why Spider is offset;
\item why our Scala and error-detection TS exceed the paper's figures;
\item whether repair's base gain is distillation;
\item why rounds do not help Instruct;
\item how much of repair-TS-R's NL-to-R gain is data volume (no same-data control for R);
\item whether repair-ED's round-1 gain comes from the feedback or from the second sample.
\end{itemize}
\end{keypoint}
